\section{第七条线：从模拟图样走向实验和卫星的可区分证据}\label{story:diagnostic}
上述研究最终都要回答：哪些观测能区分两个给出相近重联率的模型？单个四极图样曾是两流体解耦的重要标志，但它的幅度、极性、尺度及嵌套结构还受压力、多个X点和多物种影响。L213与L283尤其说明，形状相似的内外四极可来自不同过程；必须同时检查电子/离子流、时间尺度和穿越轨迹。

实验链也需要分层。先把Ren2005、GuideExp2012等直接Hall测量与理论场结构相连，再把Yoo2013实际测得的电势、光谱和离子流，与L272后续提出的随机加热解释相连。这个次序把“观察事实”和“解释候选”区分开，也为第5章的仪器尺度、积分效应与误差分析提供明确的问题。

\subsection{四极Hall场的理论结构怎样接受实验检验}
\citet{L089}从稳态不可压EMHD出发，用单位磁通的体积积分解释四极Hall场与电子流函数的关系。
该推导最有价值的细节是：理想EMHD内区的面外场主要由电子输运，生成需联系离子过渡区；预设X形面内场还在分离线产生对数奇异，必须靠惯性等修正。
它解释四极结构，未独立确定重联率和电子层长。\citet{L119}的实验只在远离X点和分离线的部分区域支持其沿射线近似恒定的场幅，近区明显偏离。

\subsection{嵌套四极：多个X点并不自动产生同样的图样}\label{sec:synthesis-nested}
\citet{L213}在EMHD中显示，强主X点的出流与弱次级X点的入流相互作用，可以形成一个不直接围绕单个X点的新外四极，包围原来的扩展四极。特别值得保留的对照是：三个同强度重联点并未产生相同嵌套结构。因此，多X点只是条件之一，\textbf{相对强度和相互作用顺序}也决定图样。

其与Cluster的比较依赖虚拟轨道：依次核对三个磁分量的过零、Hall峰以及法向电场的投影，而不是直接测到一张完整二维四极图。模拟片半厚约$7d_e$、观测约$3$--$5d_e$；法向场仅约1 nT，坐标方向不确定度可达20度，单航天器又不能唯一分离时间和空间变化。报告因此把它列为与EMHD解释相容的观测，保留替代结构和坐标误差。

这也给下一小节的尘埃嵌套四极提供一个必要对照：两者可以形似，但前者来自电子尺度多X点相互作用，后者含更慢的第三物种响应。要判别机制，应比较内外尺度比、建立时间和分物种电流，不能从“嵌套”二字直接识别模型。

\subsection{多组分带来新的Hall图样和电场抵消}
冷离子MMS事件\citep{L227}表明，不同离子群可在同一分离线具有不同磁化程度；冷离子仍随电场漂移，其对流项可与Hall项同量级。负尘埃PIC\citep{L283}进一步出现内外极性相反的嵌套四极Hall场，外侧结构随着缓慢尘埃响应而建立。后者是三物种结果，不能列成标准电子--离子两流体已解释的现象；尘埃质量比、电荷和密度只取少数组合，也没有对应实验确认。

\subsection{MRX先测到加速与下游升温，热化机制仍待区分}\label{sec:synthesis-mrx}
\citet{Yoo2013}直接测得MRX的面内电势井在下游变深、变宽，面内电场约500--800 V/m，离子跨越分离线后可加速至约$0.5\VA$。沿下游方向到约9 cm后，流速开始趋于饱和而离子温度上升；温度横向剖面比喷流更窄。图2电势、图3离子流与电功、图4光谱和温度必须连读，才构成“先加速、再下游升温”的测量链。

这一实验约束了能量进入离子的路径，却没有唯一选出随机加热。原文明确保留再磁化、下游碰撞、波粒作用、势阱反弹及装置边界等解释；下游离子均程降至2--5 cm，也不支持把全区域一概称为严格无碰撞。电子压力梯度在电势重建中不可忽略，不能把测得的整个面内电场都机械等同于单独的$\J\times\B/(en)$。

\citet{L272}后来用电场梯度估计解释同类实验，是对这些事实提出的机制候选。与本报告已核读的非对称能量收支实验\citep{Yamada2018}相连后，研究问题变为：不同闭合能否同时预测电势、喷流饱和位置、温升位置和分物种能流？仅使一个电场阈值超过一，还不足以完成这四项检验。

\paragraph{这一线索走到了哪里。}有区分力的实验不只寻找某个图样是否出现，而是同时约束它的尺度、相位、速度、压力和能量预算。多点原位测量或与真实诊断一致的合成观测，可以排除单轨道、视线积分和坐标误差产生的歧义。本章给出的每条机制线，均应回到这类可测的联合约束；第4章据此列出尚待解决的问题，第5章展开实验依据。
